% ===================================================================
% CANONICAL EQUATION SET: CAOS_LDA_HSI manuscript series (P3/P4/P5)
% ===================================================================
% Single source of truth for all display equations. Authored once,
% \input into P3 (journal_v_sweep), P4 (journal_backbone_factorial),
% P5 (journal_interpretability), and transcribed (KaTeX) into the
% web-app methodology route and the wiki Mathematical-Background page.
%
% Usage:  % ===================================================================
% CANONICAL EQUATION SET: CAOS_LDA_HSI manuscript series (P3/P4/P5)
% ===================================================================
% Single source of truth for all display equations. Authored once,
% \input into P3 (journal_v_sweep), P4 (journal_backbone_factorial),
% P5 (journal_interpretability), and transcribed (KaTeX) into the
% web-app methodology route and the wiki Mathematical-Background page.
%
% Usage:  % ===================================================================
% CANONICAL EQUATION SET: CAOS_LDA_HSI manuscript series (P3/P4/P5)
% ===================================================================
% Single source of truth for all display equations. Authored once,
% \input into P3 (journal_v_sweep), P4 (journal_backbone_factorial),
% P5 (journal_interpretability), and transcribed (KaTeX) into the
% web-app methodology route and the wiki Mathematical-Background page.
%
% Usage:  % ===================================================================
% CANONICAL EQUATION SET: CAOS_LDA_HSI manuscript series (P3/P4/P5)
% ===================================================================
% Single source of truth for all display equations. Authored once,
% \input into P3 (journal_v_sweep), P4 (journal_backbone_factorial),
% P5 (journal_interpretability), and transcribed (KaTeX) into the
% web-app methodology route and the wiki Mathematical-Background page.
%
% Usage:  \input{../../equations/canonical.tex}  then reference the
% groups by \eqref{eq:lda-gen}, \eqref{eq:etm-beta}, etc. Groups can
% also be copied verbatim into a single \begin{align} if a paper wants
% them inline. Each equation is independently labelled.
%
% NOTATION (fixed once, used everywhere):
%   b in {1..B}    spectral band            (B ~ 200)
%   q in {0..Q-1}  quantisation level       (Q in {8,16,32}, default 8)
%   k in {1..K}    topic                    (K per-scene K-policy)
%   d in {1..D}    document = labelled pixel
%   v in V         wordification vocabulary (recipe-specific)
%   x_d in R^B     normalised pixel spectrum
%   y_d            land-cover / mineral label of pixel d
%   theta_d        per-doc topic proportions  (simplex Delta^{K-1})
%   phi_k          topic-word distribution    (simplex Delta^{|V|-1})
%   n_d in N^{|V|} doc-term count vector produced by a recipe
% ===================================================================

% ----- GROUP A : the four topic-model backbones -------------------

% A1  LDA (online variational Bayes; P1/P3 baseline)
\begin{equation}\label{eq:lda-gen}
\begin{gathered}
\phi_k \sim \mathrm{Dir}(\beta),\quad
\theta_d \sim \mathrm{Dir}(\alpha),\\
z_{dn} \sim \mathrm{Cat}(\theta_d),\quad
w_{dn} \sim \mathrm{Cat}(\phi_{z_{dn}}).
\end{gathered}
\end{equation}

% A2  HDP (nonparametric K via two-level stick-breaking)
\begin{equation}\label{eq:hdp-gen}
\begin{gathered}
\boldsymbol{\beta}\sim\mathrm{GEM}(\gamma),\quad
\pi_d \sim \mathrm{DP}(\alpha_0,\boldsymbol{\beta}),\\
z_{dn}\sim\mathrm{Cat}(\pi_d),\quad
w_{dn}\sim\mathrm{Cat}(\phi_{z_{dn}}),
\end{gathered}
\end{equation}
with global atom weights $\beta_k=\beta'_k\prod_{l<k}(1-\beta'_l)$,
$\beta'_k\sim\mathrm{Beta}(1,\gamma)$.

% A3  "ProdLDA" as run: the amortised neural topic model of Srivastava and Sutton
%     in its mixture-decoder form (their NVLDA), with a standard normal prior on
%     the topic-proportion logits delta_d and a free matrix of unnormalised
%     topic-word logits omega_k (no product of experts, no Laplace approximation)
\begin{equation}\label{eq:prodlda-gen}
\begin{gathered}
\delta_d \sim \mathcal{N}(0, I),\quad
\theta_d = \mathrm{softmax}(\delta_d),\\
\beta_k = \mathrm{softmax}(\omega_k),\quad
w_{dn} \sim \mathrm{Cat}\!\Big(\textstyle\sum_k \theta_{dk}\,\beta_k\Big),
\end{gathered}
\end{equation}
where the unnormalised topic--word logits $\omega_k\in\mathbb{R}^{|V|}$ form
a free $K\times|V|$ matrix and inference is amortised,
$(\mu_d,\Sigma_d)=\mathrm{Enc}_\eta(n_d)$. ProdLDA as published mixes the
logits instead, $w_{dn}\sim\mathrm{Cat}\big(\mathrm{softmax}(\textstyle\sum_k \theta_{dk}\,\omega_k)\big)$
(a product of experts), with a Laplace approximation to $\mathrm{Dir}(\alpha)$
as the prior; that form was not run.

% A4  ETM (embedded topic model; same prior and encoder as A3, rank-L decoder)
\begin{equation}\label{eq:etm-beta}
\begin{gathered}
\delta_d \sim \mathcal{N}(0, I),\quad
\theta_d = \mathrm{softmax}(\delta_d),\\
\beta_k = \mathrm{softmax}\!\big(\rho^{\top}\alpha_k\big),\quad
w_{dn}\sim\mathrm{Cat}\!\Big(\textstyle\sum_k \theta_{dk}\,\beta_k\Big),
\end{gathered}
\end{equation}
with word embeddings $\rho\in\mathbb{R}^{L\times|V|}$ and topic
embeddings $\alpha_k\in\mathbb{R}^{L}$ ($L=64$), and the same amortised
encoder as \eqref{eq:prodlda-gen}.

% ----- GROUP B : wordification recipes (spectrum -> doc-term) ------

% B0  per-row min-max normalisation + uniform quantiser
\begin{equation}\label{eq:quantiser}
\begin{gathered}
\tilde{x}_{db}=\frac{x_{db}-\min_b x_{db}}{\max_b x_{db}-\min_b x_{db}},
\\
q_b(x_d)=\mathrm{clip}\big(\lfloor \tilde{x}_{db}\,Q\rfloor,\,0,\,Q-1\big).
\end{gathered}
\end{equation}

% B1  V1 band-frequency: token = band, multiplicity = quantised intensity
\begin{equation}\label{eq:v1}
n_d^{\mathrm{V1}}[b] = q_b(x_d),\qquad |V_{\mathrm{V1}}| = B.
\end{equation}

% B2  V3 joint (band, bin): one token per band naming its bin
\begin{equation}\label{eq:v3}
n_d^{\mathrm{V3}}[(b,q)] = \mathds{1}\!\left[q_b(x_d)=q\right],\qquad
|V_{\mathrm{V3}}| = B\,Q .
\end{equation}

% B3  V8 NFINDR endmember-fraction: FCLS abundances, rescaled by the pixel's largest abundance and
%     quantised to Q levels (build_wordifications_v6plus.wordify_v8_endmember_fraction)
\begin{equation}\label{eq:v8}
\begin{gathered}
a_d=\arg\min_{a\ge 0,\;\mathbf{1}^{\top}a=1}\;\lVert x_d - E\,a\rVert_2^2,\\
n_d^{\mathrm{V8}}[m]=\min\!\Big(\Big\lfloor Q\,\frac{a_{dm}}{\max_{m'} a_{dm'}}\Big\rfloor,\;Q-1\Big),
\end{gathered}
\end{equation}
where $E=[e_1,\dots,e_M]$ are the NFINDR endmembers (simplex vertices),
$x_d$ is the unnormalised spectrum and $|V_{\mathrm{V8}}|=M$. The abundance
non-negativity (ANC, $a\ge 0$) and sum-to-one (ASC, $\mathbf{1}^{\top}a=1$)
constraints make $a_d$ a point on the simplex spanned by the endmembers
(fully-constrained least squares, FCLS; implemented as
$\delta$-augmented NNLS followed by renormalisation), so the coefficients
are proper abundance fractions over the convex hull. The dominant
endmember of every pixel emits $Q-1$ tokens, so the document length
grows with $Q$.

% B4  V12 GMM-token: (band, component) joints from one 1-D Q-component Gaussian mixture fitted to the
%     band intensities (build_wordifications_v6plus.wordify_v12_gmm; hard assignment, not responsibilities)
\begin{equation}\label{eq:v12}
\begin{gathered}
g_b(x_d)=\arg\max_{k\in\{0,\dots,Q-1\}}\;\pi_k\,\mathcal{N}\big(x_{db}\mid\mu_k,\sigma_k^2\big),\\
n_d^{\mathrm{V12}}[(b,g)]=\mathds{1}\!\left[g_b(x_d)=g\right],\qquad
|V_{\mathrm{V12}}|=B\,Q,
\end{gathered}
\end{equation}
where $(\pi_k,\mu_k,\sigma_k^2)_{k<Q}$ is a one-dimensional Gaussian
mixture fitted once per scene to the unnormalised band intensities of the
sampled pixels (a random subset of 50\,000 values), with the components
indexed in increasing order of $\mu_k$. Each band emits one token, as in
V3, but the bins are shared by all pixels and follow the intensity
distribution instead of splitting each spectrum's own range into equal
intervals.

% B5  V20 MI-weighted bands: V3 alphabet, per-band emission multiplicity
\begin{equation}\label{eq:v20}
\begin{gathered}
w_b=\Big\lfloor \tfrac{\mathrm{MI}(x_{\cdot b};\,y)}{\max_{b'}\mathrm{MI}(x_{\cdot b'};\,y)}\;Q_{\max}\Big\rceil,
\\
n_d^{\mathrm{V20}}[(b,q)] = w_b\,\mathds{1}\!\left[q_b(x_d)=q\right].
\end{gathered}
\end{equation}
The token alphabet is V3's, $|V_{\mathrm{V20}}|=B\,Q$ (=1600 at
$B=200,Q=8$; mean $\approx 1376$ across the six labelled scenes since
$B$ varies). Bands with $w_b=0$ emit nothing, so the \emph{effective}
(non-empty-column) vocabulary is smaller still ($\approx 770$ mean).

% ----- GROUP C : evaluation axes ----------------------------------

% C1  F-1 macro-F1 (5-fold; topic_routed_soft head)
\begin{equation}\label{eq:f1}
\mathrm{F\text{-}1}=\frac{1}{|\mathcal{C}|}\sum_{c\in\mathcal{C}}
\frac{2\,P_c R_c}{P_c+R_c}.
\end{equation}

% C2  F-2 coherence c_v (NPMI sliding window + cosine aggregation)
\begin{equation}\label{eq:f2}
\begin{gathered}
\mathrm{NPMI}(w_i,w_j)=\frac{\log\frac{p(w_i,w_j)+\epsilon}{p(w_i)p(w_j)}}{-\log\big(p(w_i,w_j)+\epsilon\big)},\\
c_v=\frac{1}{N}\sum_{i}\cos\!\big(\mathbf{v}_i,\,\textstyle\sum_j \mathbf{v}_j\big),
\end{gathered}
\end{equation}
with $\mathbf{v}_i=\big(\mathrm{NPMI}(w_i,w_\ell)\big)_{\ell=1}^{N}$ over the
top-$N$ topic words.

% C3  F-7 topic-to-label MI normalised by H(label) -- this is the
% ASYMMETRIC label-normalised MI (uncertainty coefficient U(Y|Zhat) =
% completeness, Rosenberg-Hirschberg 2007), NOT the symmetric NMI
% I/sqrt(H H). Uncorrected for chance.
\begin{equation}\label{eq:f7}
\begin{gathered}
\hat{z}_d=\arg\max_k \theta_{dk},\\
\mathrm{F\text{-}7}=\frac{I(\hat{Z};Y)}{H(Y)}
\;=\; U(Y\mid\hat{Z})\;=\;1-\frac{H(Y\mid\hat{Z})}{H(Y)}.
\end{gathered}
\end{equation}

% C4  F-14 top-N pairwise Jaccard repetitiveness (lower = more diverse)
\begin{equation}\label{eq:f14}
\mathrm{F\text{-}14}=\frac{2}{K(K-1)}\sum_{k<k'}
\frac{|T_k\cap T_{k'}|}{|T_k\cup T_{k'}|},\qquad T_k=\text{top-}N(\phi_k).
\end{equation}

% C5  F-18 Maier reseed reliability (Hungarian-matched top-N cosine)
\begin{equation}\label{eq:f18}
\mathrm{F\text{-}18}=\frac{1}{\binom{S}{2}}\sum_{s<s'}\frac{1}{K}
\max_{\sigma\in\mathfrak{S}_K}\sum_{k}\cos\!\big(\mathbf{u}^{(s)}_k,\mathbf{u}^{(s')}_{\sigma(k)}\big),
\end{equation}
over $S$ reseeded refits, with $\mathbf{u}^{(s)}_k$ the top-$N$ indicator
vector of topic $k$ in run $s$.

% C6  F-22 counterfactual L1 (min token perturbation to flip argmax topic)
\begin{equation}\label{eq:f22}
\begin{aligned}
\mathrm{F\text{-}22}(d)={}&\min_{\delta\in\mathbb{Z}^{|V|}}\lVert\delta\rVert_1
\;\;\text{s.t.}\\
&\arg\max_k\theta_k(n_d+\delta)\neq\arg\max_k\theta_k(n_d),\\
&n_d+\delta\ge 0,
\end{aligned}
\end{equation}
reported as the per-scene median, capped at a sentinel
$\mathrm{MAX\_STEPS}=50$ when no flip is found.

% C7  F-15 LLM-judge alignment (top-N overlap rule, as the judge implementation computes it:
%     ambiguous documents, i.e. neither aligned nor with >= 3 non-zero tokens, are excluded)
\begin{equation}\label{eq:f15}
\begin{aligned}
\mathrm{aligned}(d)&=\mathds{1}\big[\,|\mathrm{top}_N(n_d)\cap\mathrm{top}_N(\phi_{\hat{z}_d})|\ge\tau\\
&\qquad\;\vee\;\mathrm{top}_3(n_d)\cap\mathrm{top}_N(\phi_{\hat{z}_d})\neq\emptyset\,\big],\\
\mathrm{decisive}(d)&=\mathds{1}\big[\,\mathrm{aligned}(d)=1\;\vee\;|\mathrm{supp}(n_d)|\ge 3\,\big],\\
\mathrm{F\text{-}15}&=\textstyle\sum_d \mathrm{aligned}(d)\,\big/\textstyle\sum_d \mathrm{decisive}(d).
\end{aligned}
\end{equation}

% C8  token-mass dispersion (P5 mechanism metric: effective vocabulary)
\begin{equation}\label{eq:dispersion}
\begin{gathered}
p_{dv}=\frac{n_{dv}}{\sum_{v'} n_{dv'}},\\
N_{\mathrm{eff}}(d)=\exp\!\big(H(p_d)\big)=\exp\!\Big(-\sum_v p_{dv}\log p_{dv}\Big).
\end{gathered}
\end{equation}
$N_{\mathrm{eff}}$ counts the tokens that effectively carry the mass of
$p_d$ (or of a topic's $\phi_k$): it is low when the mass is concentrated
on few tokens and at most $|V|$.

% ----- GROUP D : cross-backbone affinity score (P4) ----------------

% D1  per-recipe cross-backbone affinity (min-max normalised within backbone)
\begin{equation}\label{eq:affinity}
\mathcal{A}(V_i)=\frac{1}{|\mathcal{B}|}\sum_{m\in\mathcal{B}}
\frac{c_v^{(m)}(V_i)-\min_j c_v^{(m)}(V_j)}{\max_j c_v^{(m)}(V_j)-\min_j c_v^{(m)}(V_j)}\;\in[0,1],
\end{equation}
averaged over backbones $\mathcal{B}=\{\mathrm{LDA},\mathrm{HDP},\mathrm{ProdLDA},\mathrm{ETM}\}$;
$\mathcal{A}\to 1$ when a recipe sits near the per-backbone top everywhere.
  then reference the
% groups by \eqref{eq:lda-gen}, \eqref{eq:etm-beta}, etc. Groups can
% also be copied verbatim into a single \begin{align} if a paper wants
% them inline. Each equation is independently labelled.
%
% NOTATION (fixed once, used everywhere):
%   b in {1..B}    spectral band            (B ~ 200)
%   q in {0..Q-1}  quantisation level       (Q in {8,16,32}, default 8)
%   k in {1..K}    topic                    (K per-scene K-policy)
%   d in {1..D}    document = labelled pixel
%   v in V         wordification vocabulary (recipe-specific)
%   x_d in R^B     normalised pixel spectrum
%   y_d            land-cover / mineral label of pixel d
%   theta_d        per-doc topic proportions  (simplex Delta^{K-1})
%   phi_k          topic-word distribution    (simplex Delta^{|V|-1})
%   n_d in N^{|V|} doc-term count vector produced by a recipe
% ===================================================================

% ----- GROUP A : the four topic-model backbones -------------------

% A1  LDA (online variational Bayes; P1/P3 baseline)
\begin{equation}\label{eq:lda-gen}
\begin{gathered}
\phi_k \sim \mathrm{Dir}(\beta),\quad
\theta_d \sim \mathrm{Dir}(\alpha),\\
z_{dn} \sim \mathrm{Cat}(\theta_d),\quad
w_{dn} \sim \mathrm{Cat}(\phi_{z_{dn}}).
\end{gathered}
\end{equation}

% A2  HDP (nonparametric K via two-level stick-breaking)
\begin{equation}\label{eq:hdp-gen}
\begin{gathered}
\boldsymbol{\beta}\sim\mathrm{GEM}(\gamma),\quad
\pi_d \sim \mathrm{DP}(\alpha_0,\boldsymbol{\beta}),\\
z_{dn}\sim\mathrm{Cat}(\pi_d),\quad
w_{dn}\sim\mathrm{Cat}(\phi_{z_{dn}}),
\end{gathered}
\end{equation}
with global atom weights $\beta_k=\beta'_k\prod_{l<k}(1-\beta'_l)$,
$\beta'_k\sim\mathrm{Beta}(1,\gamma)$.

% A3  "ProdLDA" as run: the amortised neural topic model of Srivastava and Sutton
%     in its mixture-decoder form (their NVLDA), with a standard normal prior on
%     the topic-proportion logits delta_d and a free matrix of unnormalised
%     topic-word logits omega_k (no product of experts, no Laplace approximation)
\begin{equation}\label{eq:prodlda-gen}
\begin{gathered}
\delta_d \sim \mathcal{N}(0, I),\quad
\theta_d = \mathrm{softmax}(\delta_d),\\
\beta_k = \mathrm{softmax}(\omega_k),\quad
w_{dn} \sim \mathrm{Cat}\!\Big(\textstyle\sum_k \theta_{dk}\,\beta_k\Big),
\end{gathered}
\end{equation}
where the unnormalised topic--word logits $\omega_k\in\mathbb{R}^{|V|}$ form
a free $K\times|V|$ matrix and inference is amortised,
$(\mu_d,\Sigma_d)=\mathrm{Enc}_\eta(n_d)$. ProdLDA as published mixes the
logits instead, $w_{dn}\sim\mathrm{Cat}\big(\mathrm{softmax}(\textstyle\sum_k \theta_{dk}\,\omega_k)\big)$
(a product of experts), with a Laplace approximation to $\mathrm{Dir}(\alpha)$
as the prior; that form was not run.

% A4  ETM (embedded topic model; same prior and encoder as A3, rank-L decoder)
\begin{equation}\label{eq:etm-beta}
\begin{gathered}
\delta_d \sim \mathcal{N}(0, I),\quad
\theta_d = \mathrm{softmax}(\delta_d),\\
\beta_k = \mathrm{softmax}\!\big(\rho^{\top}\alpha_k\big),\quad
w_{dn}\sim\mathrm{Cat}\!\Big(\textstyle\sum_k \theta_{dk}\,\beta_k\Big),
\end{gathered}
\end{equation}
with word embeddings $\rho\in\mathbb{R}^{L\times|V|}$ and topic
embeddings $\alpha_k\in\mathbb{R}^{L}$ ($L=64$), and the same amortised
encoder as \eqref{eq:prodlda-gen}.

% ----- GROUP B : wordification recipes (spectrum -> doc-term) ------

% B0  per-row min-max normalisation + uniform quantiser
\begin{equation}\label{eq:quantiser}
\begin{gathered}
\tilde{x}_{db}=\frac{x_{db}-\min_b x_{db}}{\max_b x_{db}-\min_b x_{db}},
\\
q_b(x_d)=\mathrm{clip}\big(\lfloor \tilde{x}_{db}\,Q\rfloor,\,0,\,Q-1\big).
\end{gathered}
\end{equation}

% B1  V1 band-frequency: token = band, multiplicity = quantised intensity
\begin{equation}\label{eq:v1}
n_d^{\mathrm{V1}}[b] = q_b(x_d),\qquad |V_{\mathrm{V1}}| = B.
\end{equation}

% B2  V3 joint (band, bin): one token per band naming its bin
\begin{equation}\label{eq:v3}
n_d^{\mathrm{V3}}[(b,q)] = \mathds{1}\!\left[q_b(x_d)=q\right],\qquad
|V_{\mathrm{V3}}| = B\,Q .
\end{equation}

% B3  V8 NFINDR endmember-fraction: FCLS abundances, rescaled by the pixel's largest abundance and
%     quantised to Q levels (build_wordifications_v6plus.wordify_v8_endmember_fraction)
\begin{equation}\label{eq:v8}
\begin{gathered}
a_d=\arg\min_{a\ge 0,\;\mathbf{1}^{\top}a=1}\;\lVert x_d - E\,a\rVert_2^2,\\
n_d^{\mathrm{V8}}[m]=\min\!\Big(\Big\lfloor Q\,\frac{a_{dm}}{\max_{m'} a_{dm'}}\Big\rfloor,\;Q-1\Big),
\end{gathered}
\end{equation}
where $E=[e_1,\dots,e_M]$ are the NFINDR endmembers (simplex vertices),
$x_d$ is the unnormalised spectrum and $|V_{\mathrm{V8}}|=M$. The abundance
non-negativity (ANC, $a\ge 0$) and sum-to-one (ASC, $\mathbf{1}^{\top}a=1$)
constraints make $a_d$ a point on the simplex spanned by the endmembers
(fully-constrained least squares, FCLS; implemented as
$\delta$-augmented NNLS followed by renormalisation), so the coefficients
are proper abundance fractions over the convex hull. The dominant
endmember of every pixel emits $Q-1$ tokens, so the document length
grows with $Q$.

% B4  V12 GMM-token: (band, component) joints from one 1-D Q-component Gaussian mixture fitted to the
%     band intensities (build_wordifications_v6plus.wordify_v12_gmm; hard assignment, not responsibilities)
\begin{equation}\label{eq:v12}
\begin{gathered}
g_b(x_d)=\arg\max_{k\in\{0,\dots,Q-1\}}\;\pi_k\,\mathcal{N}\big(x_{db}\mid\mu_k,\sigma_k^2\big),\\
n_d^{\mathrm{V12}}[(b,g)]=\mathds{1}\!\left[g_b(x_d)=g\right],\qquad
|V_{\mathrm{V12}}|=B\,Q,
\end{gathered}
\end{equation}
where $(\pi_k,\mu_k,\sigma_k^2)_{k<Q}$ is a one-dimensional Gaussian
mixture fitted once per scene to the unnormalised band intensities of the
sampled pixels (a random subset of 50\,000 values), with the components
indexed in increasing order of $\mu_k$. Each band emits one token, as in
V3, but the bins are shared by all pixels and follow the intensity
distribution instead of splitting each spectrum's own range into equal
intervals.

% B5  V20 MI-weighted bands: V3 alphabet, per-band emission multiplicity
\begin{equation}\label{eq:v20}
\begin{gathered}
w_b=\Big\lfloor \tfrac{\mathrm{MI}(x_{\cdot b};\,y)}{\max_{b'}\mathrm{MI}(x_{\cdot b'};\,y)}\;Q_{\max}\Big\rceil,
\\
n_d^{\mathrm{V20}}[(b,q)] = w_b\,\mathds{1}\!\left[q_b(x_d)=q\right].
\end{gathered}
\end{equation}
The token alphabet is V3's, $|V_{\mathrm{V20}}|=B\,Q$ (=1600 at
$B=200,Q=8$; mean $\approx 1376$ across the six labelled scenes since
$B$ varies). Bands with $w_b=0$ emit nothing, so the \emph{effective}
(non-empty-column) vocabulary is smaller still ($\approx 770$ mean).

% ----- GROUP C : evaluation axes ----------------------------------

% C1  F-1 macro-F1 (5-fold; topic_routed_soft head)
\begin{equation}\label{eq:f1}
\mathrm{F\text{-}1}=\frac{1}{|\mathcal{C}|}\sum_{c\in\mathcal{C}}
\frac{2\,P_c R_c}{P_c+R_c}.
\end{equation}

% C2  F-2 coherence c_v (NPMI sliding window + cosine aggregation)
\begin{equation}\label{eq:f2}
\begin{gathered}
\mathrm{NPMI}(w_i,w_j)=\frac{\log\frac{p(w_i,w_j)+\epsilon}{p(w_i)p(w_j)}}{-\log\big(p(w_i,w_j)+\epsilon\big)},\\
c_v=\frac{1}{N}\sum_{i}\cos\!\big(\mathbf{v}_i,\,\textstyle\sum_j \mathbf{v}_j\big),
\end{gathered}
\end{equation}
with $\mathbf{v}_i=\big(\mathrm{NPMI}(w_i,w_\ell)\big)_{\ell=1}^{N}$ over the
top-$N$ topic words.

% C3  F-7 topic-to-label MI normalised by H(label) -- this is the
% ASYMMETRIC label-normalised MI (uncertainty coefficient U(Y|Zhat) =
% completeness, Rosenberg-Hirschberg 2007), NOT the symmetric NMI
% I/sqrt(H H). Uncorrected for chance.
\begin{equation}\label{eq:f7}
\begin{gathered}
\hat{z}_d=\arg\max_k \theta_{dk},\\
\mathrm{F\text{-}7}=\frac{I(\hat{Z};Y)}{H(Y)}
\;=\; U(Y\mid\hat{Z})\;=\;1-\frac{H(Y\mid\hat{Z})}{H(Y)}.
\end{gathered}
\end{equation}

% C4  F-14 top-N pairwise Jaccard repetitiveness (lower = more diverse)
\begin{equation}\label{eq:f14}
\mathrm{F\text{-}14}=\frac{2}{K(K-1)}\sum_{k<k'}
\frac{|T_k\cap T_{k'}|}{|T_k\cup T_{k'}|},\qquad T_k=\text{top-}N(\phi_k).
\end{equation}

% C5  F-18 Maier reseed reliability (Hungarian-matched top-N cosine)
\begin{equation}\label{eq:f18}
\mathrm{F\text{-}18}=\frac{1}{\binom{S}{2}}\sum_{s<s'}\frac{1}{K}
\max_{\sigma\in\mathfrak{S}_K}\sum_{k}\cos\!\big(\mathbf{u}^{(s)}_k,\mathbf{u}^{(s')}_{\sigma(k)}\big),
\end{equation}
over $S$ reseeded refits, with $\mathbf{u}^{(s)}_k$ the top-$N$ indicator
vector of topic $k$ in run $s$.

% C6  F-22 counterfactual L1 (min token perturbation to flip argmax topic)
\begin{equation}\label{eq:f22}
\begin{aligned}
\mathrm{F\text{-}22}(d)={}&\min_{\delta\in\mathbb{Z}^{|V|}}\lVert\delta\rVert_1
\;\;\text{s.t.}\\
&\arg\max_k\theta_k(n_d+\delta)\neq\arg\max_k\theta_k(n_d),\\
&n_d+\delta\ge 0,
\end{aligned}
\end{equation}
reported as the per-scene median, capped at a sentinel
$\mathrm{MAX\_STEPS}=50$ when no flip is found.

% C7  F-15 LLM-judge alignment (top-N overlap rule, as the judge implementation computes it:
%     ambiguous documents, i.e. neither aligned nor with >= 3 non-zero tokens, are excluded)
\begin{equation}\label{eq:f15}
\begin{aligned}
\mathrm{aligned}(d)&=\mathds{1}\big[\,|\mathrm{top}_N(n_d)\cap\mathrm{top}_N(\phi_{\hat{z}_d})|\ge\tau\\
&\qquad\;\vee\;\mathrm{top}_3(n_d)\cap\mathrm{top}_N(\phi_{\hat{z}_d})\neq\emptyset\,\big],\\
\mathrm{decisive}(d)&=\mathds{1}\big[\,\mathrm{aligned}(d)=1\;\vee\;|\mathrm{supp}(n_d)|\ge 3\,\big],\\
\mathrm{F\text{-}15}&=\textstyle\sum_d \mathrm{aligned}(d)\,\big/\textstyle\sum_d \mathrm{decisive}(d).
\end{aligned}
\end{equation}

% C8  token-mass dispersion (P5 mechanism metric: effective vocabulary)
\begin{equation}\label{eq:dispersion}
\begin{gathered}
p_{dv}=\frac{n_{dv}}{\sum_{v'} n_{dv'}},\\
N_{\mathrm{eff}}(d)=\exp\!\big(H(p_d)\big)=\exp\!\Big(-\sum_v p_{dv}\log p_{dv}\Big).
\end{gathered}
\end{equation}
$N_{\mathrm{eff}}$ counts the tokens that effectively carry the mass of
$p_d$ (or of a topic's $\phi_k$): it is low when the mass is concentrated
on few tokens and at most $|V|$.

% ----- GROUP D : cross-backbone affinity score (P4) ----------------

% D1  per-recipe cross-backbone affinity (min-max normalised within backbone)
\begin{equation}\label{eq:affinity}
\mathcal{A}(V_i)=\frac{1}{|\mathcal{B}|}\sum_{m\in\mathcal{B}}
\frac{c_v^{(m)}(V_i)-\min_j c_v^{(m)}(V_j)}{\max_j c_v^{(m)}(V_j)-\min_j c_v^{(m)}(V_j)}\;\in[0,1],
\end{equation}
averaged over backbones $\mathcal{B}=\{\mathrm{LDA},\mathrm{HDP},\mathrm{ProdLDA},\mathrm{ETM}\}$;
$\mathcal{A}\to 1$ when a recipe sits near the per-backbone top everywhere.
  then reference the
% groups by \eqref{eq:lda-gen}, \eqref{eq:etm-beta}, etc. Groups can
% also be copied verbatim into a single \begin{align} if a paper wants
% them inline. Each equation is independently labelled.
%
% NOTATION (fixed once, used everywhere):
%   b in {1..B}    spectral band            (B ~ 200)
%   q in {0..Q-1}  quantisation level       (Q in {8,16,32}, default 8)
%   k in {1..K}    topic                    (K per-scene K-policy)
%   d in {1..D}    document = labelled pixel
%   v in V         wordification vocabulary (recipe-specific)
%   x_d in R^B     normalised pixel spectrum
%   y_d            land-cover / mineral label of pixel d
%   theta_d        per-doc topic proportions  (simplex Delta^{K-1})
%   phi_k          topic-word distribution    (simplex Delta^{|V|-1})
%   n_d in N^{|V|} doc-term count vector produced by a recipe
% ===================================================================

% ----- GROUP A : the four topic-model backbones -------------------

% A1  LDA (online variational Bayes; P1/P3 baseline)
\begin{equation}\label{eq:lda-gen}
\begin{gathered}
\phi_k \sim \mathrm{Dir}(\beta),\quad
\theta_d \sim \mathrm{Dir}(\alpha),\\
z_{dn} \sim \mathrm{Cat}(\theta_d),\quad
w_{dn} \sim \mathrm{Cat}(\phi_{z_{dn}}).
\end{gathered}
\end{equation}

% A2  HDP (nonparametric K via two-level stick-breaking)
\begin{equation}\label{eq:hdp-gen}
\begin{gathered}
\boldsymbol{\beta}\sim\mathrm{GEM}(\gamma),\quad
\pi_d \sim \mathrm{DP}(\alpha_0,\boldsymbol{\beta}),\\
z_{dn}\sim\mathrm{Cat}(\pi_d),\quad
w_{dn}\sim\mathrm{Cat}(\phi_{z_{dn}}),
\end{gathered}
\end{equation}
with global atom weights $\beta_k=\beta'_k\prod_{l<k}(1-\beta'_l)$,
$\beta'_k\sim\mathrm{Beta}(1,\gamma)$.

% A3  "ProdLDA" as run: the amortised neural topic model of Srivastava and Sutton
%     in its mixture-decoder form (their NVLDA), with a standard normal prior on
%     the topic-proportion logits delta_d and a free matrix of unnormalised
%     topic-word logits omega_k (no product of experts, no Laplace approximation)
\begin{equation}\label{eq:prodlda-gen}
\begin{gathered}
\delta_d \sim \mathcal{N}(0, I),\quad
\theta_d = \mathrm{softmax}(\delta_d),\\
\beta_k = \mathrm{softmax}(\omega_k),\quad
w_{dn} \sim \mathrm{Cat}\!\Big(\textstyle\sum_k \theta_{dk}\,\beta_k\Big),
\end{gathered}
\end{equation}
where the unnormalised topic--word logits $\omega_k\in\mathbb{R}^{|V|}$ form
a free $K\times|V|$ matrix and inference is amortised,
$(\mu_d,\Sigma_d)=\mathrm{Enc}_\eta(n_d)$. ProdLDA as published mixes the
logits instead, $w_{dn}\sim\mathrm{Cat}\big(\mathrm{softmax}(\textstyle\sum_k \theta_{dk}\,\omega_k)\big)$
(a product of experts), with a Laplace approximation to $\mathrm{Dir}(\alpha)$
as the prior; that form was not run.

% A4  ETM (embedded topic model; same prior and encoder as A3, rank-L decoder)
\begin{equation}\label{eq:etm-beta}
\begin{gathered}
\delta_d \sim \mathcal{N}(0, I),\quad
\theta_d = \mathrm{softmax}(\delta_d),\\
\beta_k = \mathrm{softmax}\!\big(\rho^{\top}\alpha_k\big),\quad
w_{dn}\sim\mathrm{Cat}\!\Big(\textstyle\sum_k \theta_{dk}\,\beta_k\Big),
\end{gathered}
\end{equation}
with word embeddings $\rho\in\mathbb{R}^{L\times|V|}$ and topic
embeddings $\alpha_k\in\mathbb{R}^{L}$ ($L=64$), and the same amortised
encoder as \eqref{eq:prodlda-gen}.

% ----- GROUP B : wordification recipes (spectrum -> doc-term) ------

% B0  per-row min-max normalisation + uniform quantiser
\begin{equation}\label{eq:quantiser}
\begin{gathered}
\tilde{x}_{db}=\frac{x_{db}-\min_b x_{db}}{\max_b x_{db}-\min_b x_{db}},
\\
q_b(x_d)=\mathrm{clip}\big(\lfloor \tilde{x}_{db}\,Q\rfloor,\,0,\,Q-1\big).
\end{gathered}
\end{equation}

% B1  V1 band-frequency: token = band, multiplicity = quantised intensity
\begin{equation}\label{eq:v1}
n_d^{\mathrm{V1}}[b] = q_b(x_d),\qquad |V_{\mathrm{V1}}| = B.
\end{equation}

% B2  V3 joint (band, bin): one token per band naming its bin
\begin{equation}\label{eq:v3}
n_d^{\mathrm{V3}}[(b,q)] = \mathds{1}\!\left[q_b(x_d)=q\right],\qquad
|V_{\mathrm{V3}}| = B\,Q .
\end{equation}

% B3  V8 NFINDR endmember-fraction: FCLS abundances, rescaled by the pixel's largest abundance and
%     quantised to Q levels (build_wordifications_v6plus.wordify_v8_endmember_fraction)
\begin{equation}\label{eq:v8}
\begin{gathered}
a_d=\arg\min_{a\ge 0,\;\mathbf{1}^{\top}a=1}\;\lVert x_d - E\,a\rVert_2^2,\\
n_d^{\mathrm{V8}}[m]=\min\!\Big(\Big\lfloor Q\,\frac{a_{dm}}{\max_{m'} a_{dm'}}\Big\rfloor,\;Q-1\Big),
\end{gathered}
\end{equation}
where $E=[e_1,\dots,e_M]$ are the NFINDR endmembers (simplex vertices),
$x_d$ is the unnormalised spectrum and $|V_{\mathrm{V8}}|=M$. The abundance
non-negativity (ANC, $a\ge 0$) and sum-to-one (ASC, $\mathbf{1}^{\top}a=1$)
constraints make $a_d$ a point on the simplex spanned by the endmembers
(fully-constrained least squares, FCLS; implemented as
$\delta$-augmented NNLS followed by renormalisation), so the coefficients
are proper abundance fractions over the convex hull. The dominant
endmember of every pixel emits $Q-1$ tokens, so the document length
grows with $Q$.

% B4  V12 GMM-token: (band, component) joints from one 1-D Q-component Gaussian mixture fitted to the
%     band intensities (build_wordifications_v6plus.wordify_v12_gmm; hard assignment, not responsibilities)
\begin{equation}\label{eq:v12}
\begin{gathered}
g_b(x_d)=\arg\max_{k\in\{0,\dots,Q-1\}}\;\pi_k\,\mathcal{N}\big(x_{db}\mid\mu_k,\sigma_k^2\big),\\
n_d^{\mathrm{V12}}[(b,g)]=\mathds{1}\!\left[g_b(x_d)=g\right],\qquad
|V_{\mathrm{V12}}|=B\,Q,
\end{gathered}
\end{equation}
where $(\pi_k,\mu_k,\sigma_k^2)_{k<Q}$ is a one-dimensional Gaussian
mixture fitted once per scene to the unnormalised band intensities of the
sampled pixels (a random subset of 50\,000 values), with the components
indexed in increasing order of $\mu_k$. Each band emits one token, as in
V3, but the bins are shared by all pixels and follow the intensity
distribution instead of splitting each spectrum's own range into equal
intervals.

% B5  V20 MI-weighted bands: V3 alphabet, per-band emission multiplicity
\begin{equation}\label{eq:v20}
\begin{gathered}
w_b=\Big\lfloor \tfrac{\mathrm{MI}(x_{\cdot b};\,y)}{\max_{b'}\mathrm{MI}(x_{\cdot b'};\,y)}\;Q_{\max}\Big\rceil,
\\
n_d^{\mathrm{V20}}[(b,q)] = w_b\,\mathds{1}\!\left[q_b(x_d)=q\right].
\end{gathered}
\end{equation}
The token alphabet is V3's, $|V_{\mathrm{V20}}|=B\,Q$ (=1600 at
$B=200,Q=8$; mean $\approx 1376$ across the six labelled scenes since
$B$ varies). Bands with $w_b=0$ emit nothing, so the \emph{effective}
(non-empty-column) vocabulary is smaller still ($\approx 770$ mean).

% ----- GROUP C : evaluation axes ----------------------------------

% C1  F-1 macro-F1 (5-fold; topic_routed_soft head)
\begin{equation}\label{eq:f1}
\mathrm{F\text{-}1}=\frac{1}{|\mathcal{C}|}\sum_{c\in\mathcal{C}}
\frac{2\,P_c R_c}{P_c+R_c}.
\end{equation}

% C2  F-2 coherence c_v (NPMI sliding window + cosine aggregation)
\begin{equation}\label{eq:f2}
\begin{gathered}
\mathrm{NPMI}(w_i,w_j)=\frac{\log\frac{p(w_i,w_j)+\epsilon}{p(w_i)p(w_j)}}{-\log\big(p(w_i,w_j)+\epsilon\big)},\\
c_v=\frac{1}{N}\sum_{i}\cos\!\big(\mathbf{v}_i,\,\textstyle\sum_j \mathbf{v}_j\big),
\end{gathered}
\end{equation}
with $\mathbf{v}_i=\big(\mathrm{NPMI}(w_i,w_\ell)\big)_{\ell=1}^{N}$ over the
top-$N$ topic words.

% C3  F-7 topic-to-label MI normalised by H(label) -- this is the
% ASYMMETRIC label-normalised MI (uncertainty coefficient U(Y|Zhat) =
% completeness, Rosenberg-Hirschberg 2007), NOT the symmetric NMI
% I/sqrt(H H). Uncorrected for chance.
\begin{equation}\label{eq:f7}
\begin{gathered}
\hat{z}_d=\arg\max_k \theta_{dk},\\
\mathrm{F\text{-}7}=\frac{I(\hat{Z};Y)}{H(Y)}
\;=\; U(Y\mid\hat{Z})\;=\;1-\frac{H(Y\mid\hat{Z})}{H(Y)}.
\end{gathered}
\end{equation}

% C4  F-14 top-N pairwise Jaccard repetitiveness (lower = more diverse)
\begin{equation}\label{eq:f14}
\mathrm{F\text{-}14}=\frac{2}{K(K-1)}\sum_{k<k'}
\frac{|T_k\cap T_{k'}|}{|T_k\cup T_{k'}|},\qquad T_k=\text{top-}N(\phi_k).
\end{equation}

% C5  F-18 Maier reseed reliability (Hungarian-matched top-N cosine)
\begin{equation}\label{eq:f18}
\mathrm{F\text{-}18}=\frac{1}{\binom{S}{2}}\sum_{s<s'}\frac{1}{K}
\max_{\sigma\in\mathfrak{S}_K}\sum_{k}\cos\!\big(\mathbf{u}^{(s)}_k,\mathbf{u}^{(s')}_{\sigma(k)}\big),
\end{equation}
over $S$ reseeded refits, with $\mathbf{u}^{(s)}_k$ the top-$N$ indicator
vector of topic $k$ in run $s$.

% C6  F-22 counterfactual L1 (min token perturbation to flip argmax topic)
\begin{equation}\label{eq:f22}
\begin{aligned}
\mathrm{F\text{-}22}(d)={}&\min_{\delta\in\mathbb{Z}^{|V|}}\lVert\delta\rVert_1
\;\;\text{s.t.}\\
&\arg\max_k\theta_k(n_d+\delta)\neq\arg\max_k\theta_k(n_d),\\
&n_d+\delta\ge 0,
\end{aligned}
\end{equation}
reported as the per-scene median, capped at a sentinel
$\mathrm{MAX\_STEPS}=50$ when no flip is found.

% C7  F-15 LLM-judge alignment (top-N overlap rule, as the judge implementation computes it:
%     ambiguous documents, i.e. neither aligned nor with >= 3 non-zero tokens, are excluded)
\begin{equation}\label{eq:f15}
\begin{aligned}
\mathrm{aligned}(d)&=\mathds{1}\big[\,|\mathrm{top}_N(n_d)\cap\mathrm{top}_N(\phi_{\hat{z}_d})|\ge\tau\\
&\qquad\;\vee\;\mathrm{top}_3(n_d)\cap\mathrm{top}_N(\phi_{\hat{z}_d})\neq\emptyset\,\big],\\
\mathrm{decisive}(d)&=\mathds{1}\big[\,\mathrm{aligned}(d)=1\;\vee\;|\mathrm{supp}(n_d)|\ge 3\,\big],\\
\mathrm{F\text{-}15}&=\textstyle\sum_d \mathrm{aligned}(d)\,\big/\textstyle\sum_d \mathrm{decisive}(d).
\end{aligned}
\end{equation}

% C8  token-mass dispersion (P5 mechanism metric: effective vocabulary)
\begin{equation}\label{eq:dispersion}
\begin{gathered}
p_{dv}=\frac{n_{dv}}{\sum_{v'} n_{dv'}},\\
N_{\mathrm{eff}}(d)=\exp\!\big(H(p_d)\big)=\exp\!\Big(-\sum_v p_{dv}\log p_{dv}\Big).
\end{gathered}
\end{equation}
$N_{\mathrm{eff}}$ counts the tokens that effectively carry the mass of
$p_d$ (or of a topic's $\phi_k$): it is low when the mass is concentrated
on few tokens and at most $|V|$.

% ----- GROUP D : cross-backbone affinity score (P4) ----------------

% D1  per-recipe cross-backbone affinity (min-max normalised within backbone)
\begin{equation}\label{eq:affinity}
\mathcal{A}(V_i)=\frac{1}{|\mathcal{B}|}\sum_{m\in\mathcal{B}}
\frac{c_v^{(m)}(V_i)-\min_j c_v^{(m)}(V_j)}{\max_j c_v^{(m)}(V_j)-\min_j c_v^{(m)}(V_j)}\;\in[0,1],
\end{equation}
averaged over backbones $\mathcal{B}=\{\mathrm{LDA},\mathrm{HDP},\mathrm{ProdLDA},\mathrm{ETM}\}$;
$\mathcal{A}\to 1$ when a recipe sits near the per-backbone top everywhere.
  then reference the
% groups by \eqref{eq:lda-gen}, \eqref{eq:etm-beta}, etc. Groups can
% also be copied verbatim into a single \begin{align} if a paper wants
% them inline. Each equation is independently labelled.
%
% NOTATION (fixed once, used everywhere):
%   b in {1..B}    spectral band            (B ~ 200)
%   q in {0..Q-1}  quantisation level       (Q in {8,16,32}, default 8)
%   k in {1..K}    topic                    (K per-scene K-policy)
%   d in {1..D}    document = labelled pixel
%   v in V         wordification vocabulary (recipe-specific)
%   x_d in R^B     normalised pixel spectrum
%   y_d            land-cover / mineral label of pixel d
%   theta_d        per-doc topic proportions  (simplex Delta^{K-1})
%   phi_k          topic-word distribution    (simplex Delta^{|V|-1})
%   n_d in N^{|V|} doc-term count vector produced by a recipe
% ===================================================================

% ----- GROUP A : the four topic-model backbones -------------------

% A1  LDA (online variational Bayes; P1/P3 baseline)
\begin{equation}\label{eq:lda-gen}
\begin{gathered}
\phi_k \sim \mathrm{Dir}(\beta),\quad
\theta_d \sim \mathrm{Dir}(\alpha),\\
z_{dn} \sim \mathrm{Cat}(\theta_d),\quad
w_{dn} \sim \mathrm{Cat}(\phi_{z_{dn}}).
\end{gathered}
\end{equation}

% A2  HDP (nonparametric K via two-level stick-breaking)
\begin{equation}\label{eq:hdp-gen}
\begin{gathered}
\boldsymbol{\beta}\sim\mathrm{GEM}(\gamma),\quad
\pi_d \sim \mathrm{DP}(\alpha_0,\boldsymbol{\beta}),\\
z_{dn}\sim\mathrm{Cat}(\pi_d),\quad
w_{dn}\sim\mathrm{Cat}(\phi_{z_{dn}}),
\end{gathered}
\end{equation}
with global atom weights $\beta_k=\beta'_k\prod_{l<k}(1-\beta'_l)$,
$\beta'_k\sim\mathrm{Beta}(1,\gamma)$.

% A3  "ProdLDA" as run: the amortised neural topic model of Srivastava and Sutton
%     in its mixture-decoder form (their NVLDA), with a standard normal prior on
%     the topic-proportion logits delta_d and a free matrix of unnormalised
%     topic-word logits omega_k (no product of experts, no Laplace approximation)
\begin{equation}\label{eq:prodlda-gen}
\begin{gathered}
\delta_d \sim \mathcal{N}(0, I),\quad
\theta_d = \mathrm{softmax}(\delta_d),\\
\beta_k = \mathrm{softmax}(\omega_k),\quad
w_{dn} \sim \mathrm{Cat}\!\Big(\textstyle\sum_k \theta_{dk}\,\beta_k\Big),
\end{gathered}
\end{equation}
where the unnormalised topic--word logits $\omega_k\in\mathbb{R}^{|V|}$ form
a free $K\times|V|$ matrix and inference is amortised,
$(\mu_d,\Sigma_d)=\mathrm{Enc}_\eta(n_d)$. ProdLDA as published mixes the
logits instead, $w_{dn}\sim\mathrm{Cat}\big(\mathrm{softmax}(\textstyle\sum_k \theta_{dk}\,\omega_k)\big)$
(a product of experts), with a Laplace approximation to $\mathrm{Dir}(\alpha)$
as the prior; that form was not run.

% A4  ETM (embedded topic model; same prior and encoder as A3, rank-L decoder)
\begin{equation}\label{eq:etm-beta}
\begin{gathered}
\delta_d \sim \mathcal{N}(0, I),\quad
\theta_d = \mathrm{softmax}(\delta_d),\\
\beta_k = \mathrm{softmax}\!\big(\rho^{\top}\alpha_k\big),\quad
w_{dn}\sim\mathrm{Cat}\!\Big(\textstyle\sum_k \theta_{dk}\,\beta_k\Big),
\end{gathered}
\end{equation}
with word embeddings $\rho\in\mathbb{R}^{L\times|V|}$ and topic
embeddings $\alpha_k\in\mathbb{R}^{L}$ ($L=64$), and the same amortised
encoder as \eqref{eq:prodlda-gen}.

% ----- GROUP B : wordification recipes (spectrum -> doc-term) ------

% B0  per-row min-max normalisation + uniform quantiser
\begin{equation}\label{eq:quantiser}
\begin{gathered}
\tilde{x}_{db}=\frac{x_{db}-\min_b x_{db}}{\max_b x_{db}-\min_b x_{db}},
\\
q_b(x_d)=\mathrm{clip}\big(\lfloor \tilde{x}_{db}\,Q\rfloor,\,0,\,Q-1\big).
\end{gathered}
\end{equation}

% B1  V1 band-frequency: token = band, multiplicity = quantised intensity
\begin{equation}\label{eq:v1}
n_d^{\mathrm{V1}}[b] = q_b(x_d),\qquad |V_{\mathrm{V1}}| = B.
\end{equation}

% B2  V3 joint (band, bin): one token per band naming its bin
\begin{equation}\label{eq:v3}
n_d^{\mathrm{V3}}[(b,q)] = \mathds{1}\!\left[q_b(x_d)=q\right],\qquad
|V_{\mathrm{V3}}| = B\,Q .
\end{equation}

% B3  V8 NFINDR endmember-fraction: FCLS abundances, rescaled by the pixel's largest abundance and
%     quantised to Q levels (build_wordifications_v6plus.wordify_v8_endmember_fraction)
\begin{equation}\label{eq:v8}
\begin{gathered}
a_d=\arg\min_{a\ge 0,\;\mathbf{1}^{\top}a=1}\;\lVert x_d - E\,a\rVert_2^2,\\
n_d^{\mathrm{V8}}[m]=\min\!\Big(\Big\lfloor Q\,\frac{a_{dm}}{\max_{m'} a_{dm'}}\Big\rfloor,\;Q-1\Big),
\end{gathered}
\end{equation}
where $E=[e_1,\dots,e_M]$ are the NFINDR endmembers (simplex vertices),
$x_d$ is the unnormalised spectrum and $|V_{\mathrm{V8}}|=M$. The abundance
non-negativity (ANC, $a\ge 0$) and sum-to-one (ASC, $\mathbf{1}^{\top}a=1$)
constraints make $a_d$ a point on the simplex spanned by the endmembers
(fully-constrained least squares, FCLS; implemented as
$\delta$-augmented NNLS followed by renormalisation), so the coefficients
are proper abundance fractions over the convex hull. The dominant
endmember of every pixel emits $Q-1$ tokens, so the document length
grows with $Q$.

% B4  V12 GMM-token: (band, component) joints from one 1-D Q-component Gaussian mixture fitted to the
%     band intensities (build_wordifications_v6plus.wordify_v12_gmm; hard assignment, not responsibilities)
\begin{equation}\label{eq:v12}
\begin{gathered}
g_b(x_d)=\arg\max_{k\in\{0,\dots,Q-1\}}\;\pi_k\,\mathcal{N}\big(x_{db}\mid\mu_k,\sigma_k^2\big),\\
n_d^{\mathrm{V12}}[(b,g)]=\mathds{1}\!\left[g_b(x_d)=g\right],\qquad
|V_{\mathrm{V12}}|=B\,Q,
\end{gathered}
\end{equation}
where $(\pi_k,\mu_k,\sigma_k^2)_{k<Q}$ is a one-dimensional Gaussian
mixture fitted once per scene to the unnormalised band intensities of the
sampled pixels (a random subset of 50\,000 values), with the components
indexed in increasing order of $\mu_k$. Each band emits one token, as in
V3, but the bins are shared by all pixels and follow the intensity
distribution instead of splitting each spectrum's own range into equal
intervals.

% B5  V20 MI-weighted bands: V3 alphabet, per-band emission multiplicity
\begin{equation}\label{eq:v20}
\begin{gathered}
w_b=\Big\lfloor \tfrac{\mathrm{MI}(x_{\cdot b};\,y)}{\max_{b'}\mathrm{MI}(x_{\cdot b'};\,y)}\;Q_{\max}\Big\rceil,
\\
n_d^{\mathrm{V20}}[(b,q)] = w_b\,\mathds{1}\!\left[q_b(x_d)=q\right].
\end{gathered}
\end{equation}
The token alphabet is V3's, $|V_{\mathrm{V20}}|=B\,Q$ (=1600 at
$B=200,Q=8$; mean $\approx 1376$ across the six labelled scenes since
$B$ varies). Bands with $w_b=0$ emit nothing, so the \emph{effective}
(non-empty-column) vocabulary is smaller still ($\approx 770$ mean).

% ----- GROUP C : evaluation axes ----------------------------------

% C1  F-1 macro-F1 (5-fold; topic_routed_soft head)
\begin{equation}\label{eq:f1}
\mathrm{F\text{-}1}=\frac{1}{|\mathcal{C}|}\sum_{c\in\mathcal{C}}
\frac{2\,P_c R_c}{P_c+R_c}.
\end{equation}

% C2  F-2 coherence c_v (NPMI sliding window + cosine aggregation)
\begin{equation}\label{eq:f2}
\begin{gathered}
\mathrm{NPMI}(w_i,w_j)=\frac{\log\frac{p(w_i,w_j)+\epsilon}{p(w_i)p(w_j)}}{-\log\big(p(w_i,w_j)+\epsilon\big)},\\
c_v=\frac{1}{N}\sum_{i}\cos\!\big(\mathbf{v}_i,\,\textstyle\sum_j \mathbf{v}_j\big),
\end{gathered}
\end{equation}
with $\mathbf{v}_i=\big(\mathrm{NPMI}(w_i,w_\ell)\big)_{\ell=1}^{N}$ over the
top-$N$ topic words.

% C3  F-7 topic-to-label MI normalised by H(label) -- this is the
% ASYMMETRIC label-normalised MI (uncertainty coefficient U(Y|Zhat) =
% completeness, Rosenberg-Hirschberg 2007), NOT the symmetric NMI
% I/sqrt(H H). Uncorrected for chance.
\begin{equation}\label{eq:f7}
\begin{gathered}
\hat{z}_d=\arg\max_k \theta_{dk},\\
\mathrm{F\text{-}7}=\frac{I(\hat{Z};Y)}{H(Y)}
\;=\; U(Y\mid\hat{Z})\;=\;1-\frac{H(Y\mid\hat{Z})}{H(Y)}.
\end{gathered}
\end{equation}

% C4  F-14 top-N pairwise Jaccard repetitiveness (lower = more diverse)
\begin{equation}\label{eq:f14}
\mathrm{F\text{-}14}=\frac{2}{K(K-1)}\sum_{k<k'}
\frac{|T_k\cap T_{k'}|}{|T_k\cup T_{k'}|},\qquad T_k=\text{top-}N(\phi_k).
\end{equation}

% C5  F-18 Maier reseed reliability (Hungarian-matched top-N cosine)
\begin{equation}\label{eq:f18}
\mathrm{F\text{-}18}=\frac{1}{\binom{S}{2}}\sum_{s<s'}\frac{1}{K}
\max_{\sigma\in\mathfrak{S}_K}\sum_{k}\cos\!\big(\mathbf{u}^{(s)}_k,\mathbf{u}^{(s')}_{\sigma(k)}\big),
\end{equation}
over $S$ reseeded refits, with $\mathbf{u}^{(s)}_k$ the top-$N$ indicator
vector of topic $k$ in run $s$.

% C6  F-22 counterfactual L1 (min token perturbation to flip argmax topic)
\begin{equation}\label{eq:f22}
\begin{aligned}
\mathrm{F\text{-}22}(d)={}&\min_{\delta\in\mathbb{Z}^{|V|}}\lVert\delta\rVert_1
\;\;\text{s.t.}\\
&\arg\max_k\theta_k(n_d+\delta)\neq\arg\max_k\theta_k(n_d),\\
&n_d+\delta\ge 0,
\end{aligned}
\end{equation}
reported as the per-scene median, capped at a sentinel
$\mathrm{MAX\_STEPS}=50$ when no flip is found.

% C7  F-15 LLM-judge alignment (top-N overlap rule, as the judge implementation computes it:
%     ambiguous documents, i.e. neither aligned nor with >= 3 non-zero tokens, are excluded)
\begin{equation}\label{eq:f15}
\begin{aligned}
\mathrm{aligned}(d)&=\mathds{1}\big[\,|\mathrm{top}_N(n_d)\cap\mathrm{top}_N(\phi_{\hat{z}_d})|\ge\tau\\
&\qquad\;\vee\;\mathrm{top}_3(n_d)\cap\mathrm{top}_N(\phi_{\hat{z}_d})\neq\emptyset\,\big],\\
\mathrm{decisive}(d)&=\mathds{1}\big[\,\mathrm{aligned}(d)=1\;\vee\;|\mathrm{supp}(n_d)|\ge 3\,\big],\\
\mathrm{F\text{-}15}&=\textstyle\sum_d \mathrm{aligned}(d)\,\big/\textstyle\sum_d \mathrm{decisive}(d).
\end{aligned}
\end{equation}

% C8  token-mass dispersion (P5 mechanism metric: effective vocabulary)
\begin{equation}\label{eq:dispersion}
\begin{gathered}
p_{dv}=\frac{n_{dv}}{\sum_{v'} n_{dv'}},\\
N_{\mathrm{eff}}(d)=\exp\!\big(H(p_d)\big)=\exp\!\Big(-\sum_v p_{dv}\log p_{dv}\Big).
\end{gathered}
\end{equation}
$N_{\mathrm{eff}}$ counts the tokens that effectively carry the mass of
$p_d$ (or of a topic's $\phi_k$): it is low when the mass is concentrated
on few tokens and at most $|V|$.

% ----- GROUP D : cross-backbone affinity score (P4) ----------------

% D1  per-recipe cross-backbone affinity (min-max normalised within backbone)
\begin{equation}\label{eq:affinity}
\mathcal{A}(V_i)=\frac{1}{|\mathcal{B}|}\sum_{m\in\mathcal{B}}
\frac{c_v^{(m)}(V_i)-\min_j c_v^{(m)}(V_j)}{\max_j c_v^{(m)}(V_j)-\min_j c_v^{(m)}(V_j)}\;\in[0,1],
\end{equation}
averaged over backbones $\mathcal{B}=\{\mathrm{LDA},\mathrm{HDP},\mathrm{ProdLDA},\mathrm{ETM}\}$;
$\mathcal{A}\to 1$ when a recipe sits near the per-backbone top everywhere.
